\section{Gestión de memoria virtual del sistema operativo}

El diseño de la gestión de memoria de un sistema operativo depende de tres decisiones principales:

\begin{enumerate}
  \item Utilizar o no \textit{memoria virtual}.
  \item Emplear \textit{paginación}, \textit{segmentación} o una combinación de ambas.
  \item Definir las \textit{políticas y algoritmos} de gestión de memoria.
\end{enumerate}

Las dos primeras decisiones dependen del soporte del hardware, ya que la memoria virtual requiere mecanismos de traducción de direcciones. En los sistemas modernos, la memoria virtual es una característica prácticamente universal y la paginación es la técnica predominante, por lo que la mayoría de las políticas de gestión se centran en ella.

El principal objetivo de estas políticas es \textit{reducir la tasa de fallos de página} (\textit{page faults}), ya que cada fallo implica una importante sobrecarga: seleccionar la página que será reemplazada, realizar operaciones de entrada/salida para intercambiar páginas y, mientras se espera dicha operación, efectuar un cambio de contexto para ejecutar otro proceso.

No existe una política óptima para todos los sistemas. Su rendimiento depende de múltiples factores, como el tamaño de la memoria principal, la velocidad del almacenamiento secundario, la cantidad de procesos en ejecución y el patrón de acceso a memoria de cada aplicación. Por ello, los sistemas pequeños suelen utilizar políticas de propósito general, mientras que los sistemas de gran tamaño incorporan herramientas de monitoreo y ajuste para optimizar el rendimiento según la carga de trabajo.

\subsection{Política de carga (\textit{Fetch Policy})}

La \textit{política de carga} determina \textit{cuándo una página debe incorporarse a la memoria principal}. Las dos estrategias más utilizadas son:

\begin{itemize}
  \item \textit{Paginación por demanda (\textit{Demand Paging}):} una página se carga únicamente cuando ocurre una referencia a una dirección que no se encuentra en memoria. Al iniciar un proceso suelen producirse numerosos fallos de página; posteriormente, debido al \textit{principio de localidad}, la mayoría de las referencias corresponden a páginas recientemente cargadas, reduciendo significativamente la frecuencia de fallos.

  \item \textit{Precarga (\textit{Prepaging}):} además de la página solicitada, el sistema carga páginas vecinas anticipando que serán utilizadas próximamente. Esta estrategia aprovecha que los dispositivos de almacenamiento secundario transfieren bloques contiguos con mayor eficiencia que múltiples accesos individuales. Sin embargo, pierde efectividad si las páginas adicionales nunca llegan a utilizarse.
\end{itemize}

La precarga puede realizarse al iniciar un proceso o cada vez que ocurre un fallo de página. La segunda alternativa resulta preferible porque es transparente para el programador y adapta la carga al comportamiento real del proceso.

Es importante distinguir la \textit{precarga} del \textit{intercambio de procesos} (\textit{swapping}). En el \textit{swapping}, cuando un proceso es suspendido, todas sus páginas residentes son trasladadas al almacenamiento secundario y se restauran cuando el proceso vuelve a ejecutarse. En cambio, la precarga simplemente incorpora páginas adicionales de forma preventiva para reducir futuros fallos de página.

\subsection{Política de ubicación (\textit{Placement Policy})}

La \textit{política de ubicación} determina en qué parte de la memoria principal debe colocarse una porción de un proceso.

En los sistemas de \textit{segmentación pura}, esta decisión es importante y pueden utilizarse algoritmos como \textit{first-fit}, \textit{best-fit} y otros métodos de asignación de memoria.

En cambio, en los sistemas que utilizan \textit{paginación} (pura o combinada con segmentación), la ubicación suele ser irrelevante, ya que el hardware de traducción de direcciones permite acceder a cualquier combinación página--marco con la misma eficiencia.

Una excepción importante aparece en los sistemas multiprocesador con \textit{acceso no uniforme a memoria (NUMA, \textit{Nonuniform Memory Access})}. En estos sistemas, toda la memoria es compartida y puede ser accedida por cualquier procesador, pero el tiempo de acceso depende de la distancia entre el procesador y el módulo de memoria correspondiente.

Por ello, el rendimiento mejora cuando los datos se encuentran en módulos de memoria cercanos a los procesadores que los utilizan. En consecuencia, en arquitecturas NUMA resulta deseable una \textit{estrategia automática de ubicación} que asigne las páginas al módulo de memoria que ofrezca el mejor rendimiento.

\subsection{Política de reemplazo (\textit{Replacement Policy})}

La \textit{política de reemplazo} determina qué página residente debe eliminarse de la memoria principal cuando ocurre un fallo de página y es necesario cargar una nueva página.

Es importante distinguir esta política de la \textit{gestión del conjunto residente}, que se ocupa de:
\begin{itemize}
  \item Cuántos marcos de página se asignan a cada proceso.
  \item Si el reemplazo puede realizarse únicamente sobre las páginas del proceso que produjo el fallo o sobre todas las páginas residentes del sistema.
\end{itemize}

Una vez definido el conjunto de páginas candidatas, la política de reemplazo selecciona cuál será expulsada.

El objetivo de todas las políticas es reemplazar la página con \textit{menor probabilidad de ser utilizada en el futuro cercano}. Como el futuro no puede conocerse, la mayoría de los algoritmos aprovechan el \textit{principio de localidad} y utilizan el historial reciente de referencias para estimar cuáles páginas volverán a utilizarse.

Existe un compromiso entre rendimiento y complejidad: los algoritmos más sofisticados suelen aproximar mejor el comportamiento futuro, pero requieren mayor soporte de hardware y software.

\subsubsection{Bloqueo de marcos (\textit{Frame Locking})}

Algunos marcos de memoria pueden estar \textit{bloqueados}, por lo que las páginas que contienen \textit{no pueden ser reemplazadas}. Esto ocurre con regiones críticas del sistema, como parte del núcleo del sistema operativo, estructuras de control, buffers de entrada/salida y otras áreas sensibles al tiempo.

El bloqueo se implementa mediante un \textit{bit de bloqueo} asociado a cada marco, almacenado normalmente en la tabla de marcos y, en algunos sistemas, también en la tabla de páginas.

\subsubsection{Algoritmos básicos}

La \textit{política de reemplazo} determina qué página debe expulsarse de la memoria principal cuando ocurre un fallo de página y no existen marcos libres. El objetivo es seleccionar la página con menor probabilidad de ser utilizada en el futuro cercano, aprovechando el \textit{principio de localidad}. En general, cuanto más sofisticado es el algoritmo, mayor es la sobrecarga de implementación.

Los algoritmos más importantes son:
\begin{itemize}
  \item \textit{Óptimo (Optimal):} reemplaza la página cuya próxima referencia ocurrirá más tarde. Produce el menor número posible de fallos de página, pero es imposible de implementar porque requiere conocer el futuro. Se utiliza como referencia para evaluar otros algoritmos.

  \item \textit{Least Recently Used (LRU):} reemplaza la página que lleva más tiempo sin ser referenciada. Se basa en que una página utilizada recientemente tiene mayor probabilidad de volver a utilizarse. Su rendimiento es cercano al algoritmo óptimo, pero requiere registrar continuamente el historial de accesos, lo que implica una elevada sobrecarga.

  \item \textit{First-In First-Out (FIFO):} reemplaza la página que lleva más tiempo en memoria. Es muy sencillo de implementar mediante una cola circular, pero suele ofrecer un rendimiento inferior porque puede expulsar páginas que todavía son utilizadas con frecuencia.

  \item \textit{Clock (segunda oportunidad):} aproxima el comportamiento de LRU con un costo mucho menor. Cada marco posee un \textit{bit de uso} (\textit{use bit}) que se activa cuando la página es referenciada. Los marcos se recorren de forma circular mediante un puntero:
  \begin{enumerate}
    \item Si el bit de uso es 0, la página se reemplaza.
    \item Si el bit de uso es 1, se reinicia a 0 y el algoritmo continúa con el siguiente marco.
  \end{enumerate}
  De esta forma, las páginas utilizadas recientemente reciben una ``segunda oportunidad'' antes de ser reemplazadas.
\end{itemize}

\begin{figure}[H]
  \centering
  % Insertar aquí una figura comparando los algoritmos
  % Optimal, LRU, FIFO y Clock sobre una misma secuencia de referencias.
  \caption{Comparación del comportamiento de los algoritmos de reemplazo de páginas.}
\end{figure}

Diversos estudios muestran que \textit{LRU} y \textit{Clock} presentan un rendimiento muy cercano al algoritmo óptimo, mientras que \textit{FIFO} suele generar un número considerablemente mayor de fallos de página. Debido a su buena relación entre rendimiento y costo de implementación, las variantes del algoritmo Clock son ampliamente utilizadas en sistemas operativos modernos.

\subsubsection{Algoritmo Clock mejorado}

Una mejora del algoritmo Clock incorpora, además del bit de uso (\(u\)), el \textit{bit de modificación} (\(m\)), que indica si la página fue escrita desde que se cargó en memoria. Cada página pertenece a una de las siguientes categorías:

\begin{enumerate}
  \item \(u=0,\; m=0\): no utilizada recientemente y no modificada.
  \item \(u=1,\; m=0\): utilizada recientemente y no modificada.
  \item \(u=0,\; m=1\): no utilizada recientemente y modificada.
  \item \(u=1,\; m=1\): utilizada recientemente y modificada.
\end{enumerate}

El algoritmo prioriza el reemplazo de páginas \((u=0,\; m=0)\), ya que no fueron utilizadas recientemente y, además, no necesitan escribirse nuevamente al almacenamiento secundario. Si no existen páginas de esta clase, busca páginas \((u=0,\; m=1)\); aunque requieren escritura antes de ser reemplazadas, siguen siendo mejores candidatas que las utilizadas recientemente.

Esta estrategia reduce el tiempo asociado al reemplazo de páginas y constituye una mejora práctica del algoritmo Clock simple.

\subsubsection{Búfer de páginas (\textit{Page Buffering})}

Una estrategia para mejorar el rendimiento de la memoria virtual consiste en utilizar un \textit{búfer de páginas}. Su objetivo es reducir el costo de los fallos de página y permitir el uso de algoritmos de reemplazo simples, como \textit{FIFO}.

Cuando una página es reemplazada, no se elimina inmediatamente de la memoria. En su lugar, se incorpora a una de dos listas:

\begin{itemize}
  \item \textit{Lista de páginas libres:} contiene páginas que no fueron modificadas y cuyos marcos pueden reutilizarse inmediatamente.
  \item \textit{Lista de páginas modificadas:} almacena páginas que deben escribirse en el almacenamiento secundario antes de liberar sus marcos.
\end{itemize}

Si un proceso vuelve a referenciar una página que permanece en alguna de estas listas, esta puede restaurarse rápidamente sin necesidad de leerla nuevamente desde disco. De esta forma, ambas listas funcionan como una \textit{caché de páginas}.

Además, las páginas modificadas pueden escribirse al almacenamiento secundario en bloques (\textit{clusters}) en lugar de hacerlo individualmente, reduciendo la cantidad de operaciones de entrada/salida y mejorando el rendimiento del sistema.

\begin{figure}[H]
  \centering
  % Esquema del Page Buffering:
  % Resident Set --> Reemplazo --> Lista de páginas libres (no modificadas)
  %                              --> Lista de páginas modificadas
  % Las páginas pueden volver al conjunto residente si son referenciadas.
  \caption{Funcionamiento del búfer de páginas.}
\end{figure}

\subsubsection{Política de reemplazo y tamaño de la caché}

En los sistemas modernos, el aumento del tamaño de la memoria caché hace que la política de reemplazo de páginas también influya sobre su rendimiento. Cuando una página es expulsada de la memoria principal, los bloques correspondientes almacenados en la caché también se invalidan, lo que puede incrementar la cantidad de fallos de caché.

Para minimizar este efecto, algunos sistemas complementan la política de reemplazo con una \textit{política de ubicación de páginas} en el búfer. El objetivo es distribuir las páginas consecutivas en marcos que se asignen a distintas posiciones de la caché, reduciendo conflictos y evitando que varias páginas compitan por las mismas líneas de caché.

Una estrategia adecuada de ubicación puede disminuir significativamente la cantidad de fallos de caché respecto de una asignación arbitraria.

\subsection{Gestión del conjunto residente}

El \textit{conjunto residente} (\textit{resident set}) es el conjunto de páginas de un proceso que se encuentran en memoria principal en un instante determinado. Como no es necesario cargar todas las páginas de un proceso para ejecutarlo, el sistema operativo debe decidir cuántos marcos asignar a cada proceso.

Esta decisión implica un compromiso:

\begin{itemize}
  \item Asignar \textit{pocos marcos} permite mantener más procesos en memoria y aumenta el grado de multiprogramación, pero incrementa la probabilidad de fallos de página.
  \item Asignar \textit{muchos marcos} reduce los fallos de página; sin embargo, una vez que se satisface la localidad del proceso, agregar más memoria produce escasas mejoras.
\end{itemize}

Por este motivo, existen dos políticas principales.

\subsubsection{Tamaño del conjunto residente}

\begin{itemize}
  \item \textit{Asignación fija:} cada proceso recibe un número fijo de marcos al crearse. Cuando ocurre un fallo de página, la página reemplazada debe pertenecer al mismo proceso.
  \item \textit{Asignación variable:} el número de marcos puede modificarse durante la ejecución. Un proceso con una alta tasa de fallos de página puede recibir marcos adicionales, mientras que otro con pocos fallos puede ceder parte de los suyos. Esta estrategia suele ofrecer un mejor rendimiento, aunque requiere que el sistema operativo supervise continuamente el comportamiento de los procesos.
\end{itemize}

\subsubsection{Alcance del reemplazo}

Cuando no existen marcos libres, el sistema operativo debe decidir de dónde obtener el marco que será reutilizado. Existen dos alternativas:
\begin{itemize}
  \item \textit{Reemplazo local:} la página reemplazada debe pertenecer al mismo proceso que produjo el fallo de página.
  \item \textit{Reemplazo global:} cualquier página desbloqueada de la memoria principal puede ser reemplazada, independientemente del proceso al que pertenezca.
\end{itemize}

La asignación fija implica necesariamente un \textit{reemplazo local}, ya que el tamaño del conjunto residente permanece constante. En cambio, la asignación variable puede combinarse con reemplazo local o global, permitiendo que los procesos ganen o pierdan marcos según sus necesidades.

\begin{figure}[H]
    \centering
    % Comparación entre:
    % - Asignación fija + reemplazo local.
    % - Asignación variable + reemplazo local.
    % - Asignación variable + reemplazo global.
    \caption{Relación entre el tamaño del conjunto residente y el alcance de la política de reemplazo.}
\end{figure}


\subsubsection{Asignación fija con reemplazo local}

En esta estrategia, cada proceso dispone de un \textit{número fijo de marcos} durante toda su ejecución. Cuando ocurre un fallo de página, el sistema operativo debe reemplazar una página perteneciente al mismo proceso, utilizando alguno de los algoritmos de reemplazo (FIFO, LRU, Clock, etc.).

La principal dificultad consiste en determinar el tamaño adecuado del conjunto residente:
\begin{itemize}
  \item Si se asignan \textit{muy pocos marcos}, aumentará la tasa de fallos de página y disminuirá el rendimiento del sistema.
  \item Si se asignan \textit{demasiados marcos}, se reducirá el número de procesos que pueden permanecer simultáneamente en memoria, desaprovechando el grado de multiprogramación.
\end{itemize}

\subsubsection{Asignación variable con reemplazo global}

En esta política, cada proceso puede \textit{ganar o perder marcos} durante su ejecución. Mientras existan marcos libres, un proceso que produzca un fallo de página recibe uno adicional, aumentando su conjunto residente y reduciendo la probabilidad de futuros fallos.

Cuando ya no existen marcos disponibles, el sistema operativo selecciona una página para reemplazar entre \textit{todas las páginas desbloqueadas} de la memoria principal, independientemente del proceso al que pertenezcan. Como consecuencia, un proceso puede perder un marco para beneficiar a otro.

Esta estrategia es sencilla de implementar, pero el proceso que pierde un marco no siempre es el más conveniente, lo que puede degradar el rendimiento. Una forma de reducir este problema es utilizar un \textit{búfer de páginas}, que permite recuperar rápidamente una página recientemente reemplazada si vuelve a ser referenciada.

\subsubsection{Asignación variable con reemplazo local}

En esta estrategia, cada proceso dispone inicialmente de un conjunto residente de tamaño determinado. Cuando ocurre un fallo de página, la página reemplazada se selecciona únicamente entre las páginas del mismo proceso (reemplazo local). Periódicamente, el sistema operativo evalúa el comportamiento del proceso y ajusta el tamaño de su conjunto residente, asignándole más o menos marcos según sus necesidades.

Esta política es más compleja que el reemplazo global, ya que requiere monitorear la actividad de cada proceso, pero suele ofrecer un mejor rendimiento al adaptar dinámicamente la memoria asignada.

\subsubsection{Conjunto de trabajo}

El \textit{conjunto de trabajo}, propuesto por Denning, es el conjunto de páginas referenciadas por un proceso durante las últimas \(\Delta\) unidades de tiempo virtual. Representa la localidad actual del proceso y constituye una estimación de las páginas que probablemente necesitará en el futuro cercano.

El tamaño del conjunto de trabajo depende de:
\begin{itemize}
  \item El parámetro \(\Delta\): cuanto mayor es la ventana de observación, mayor será el conjunto de trabajo.
  \item El instante de ejecución: cuando un proceso cambia de localidad, su conjunto de trabajo también cambia.
\end{itemize}

La estrategia basada en el conjunto de trabajo consiste en:
\begin{enumerate}
  \item Monitorear el conjunto de trabajo de cada proceso.
  \item Eliminar del conjunto residente las páginas que ya no pertenecen al conjunto de trabajo.
  \item Permitir la ejecución únicamente cuando el conjunto de trabajo esté completamente cargado en memoria.
\end{enumerate}

Aunque esta estrategia aprovecha eficazmente el principio de localidad, resulta difícil de implementar porque exige registrar continuamente todas las referencias a memoria y determinar un valor adecuado para \(\Delta\).

\begin{figure}[H]
  \centering
  % Figura: evolución temporal del Working Set mostrando
  % cómo cambia de tamaño al pasar de una localidad a otra.
  \caption{Variación del conjunto de trabajo durante la ejecución de un proceso.}
\end{figure}

\subsubsection{Frecuencia de fallos de página (PFF)}

Como aproximación al conjunto de trabajo, muchos sistemas monitorean la \textit{frecuencia de fallos de página} (\textit{Page Fault Frequency, PFF}).

La idea es sencilla:
\begin{itemize}
  \item Si la tasa de fallos es \textit{alta}, el proceso necesita más marcos, por lo que se incrementa su conjunto residente.
  \item Si la tasa de fallos es \textit{baja}, el proceso puede ceder parte de sus marcos sin afectar significativamente su rendimiento.
\end{itemize}

El algoritmo PFF utiliza el tiempo transcurrido entre dos fallos de página consecutivos como una estimación de la tasa de fallos. Aunque su implementación es simple, presenta una limitación importante: durante los cambios de localidad el conjunto residente puede crecer excesivamente antes de liberar las páginas que dejaron de utilizarse.

\subsubsection{Working Set Muestreado de Intervalo Variable (VSWS)}

El algoritmo \textit{Variable-Interval Sampled Working Set (VSWS)} mejora el comportamiento de PFF mediante evaluaciones periódicas del conjunto residente.

Al comienzo de cada intervalo se reinician los bits de uso de las páginas. Al finalizar el intervalo, únicamente permanecen en el conjunto residente aquellas páginas cuyo bit de uso fue activado durante dicho período; las demás son descartadas. Mientras transcurre el intervalo, las nuevas páginas referenciadas se incorporan normalmente al conjunto residente.

VSWS utiliza tres parámetros:
\begin{itemize}
  \item \(M\): duración mínima del intervalo de muestreo.
  \item \(L\): duración máxima del intervalo.
  \item \(Q\): cantidad máxima de fallos de página permitidos antes de realizar una nueva evaluación.
\end{itemize}

Cuando aumenta la frecuencia de fallos de página, el algoritmo reduce el intervalo entre evaluaciones, eliminando más rápidamente las páginas que dejaron de utilizarse. De esta forma, se adapta mejor a los cambios de localidad que el algoritmo PFF, manteniendo una complejidad de implementación similar.

\subsection{Política de vaciado/limpieza}

La \textit{política de limpieza} determina cuándo una página modificada (\textit{dirty page}) debe escribirse nuevamente al almacenamiento secundario. Las dos estrategias principales son:
\begin{itemize}
  \item \textit{Limpieza por demanda:} la página se escribe únicamente cuando va a ser reemplazada. Reduce las escrituras innecesarias, pero el proceso que produjo el fallo de página debe esperar tanto la escritura de la página modificada como la lectura de la nueva página.

  \item \textit{Limpieza anticipada:} las páginas modificadas se escriben antes de que sus marcos sean necesarios, generalmente en bloques. Esto disminuye el tiempo de espera durante un fallo de página, aunque puede generar escrituras innecesarias si la página vuelve a modificarse antes de ser reemplazada.
\end{itemize}

Una estrategia más eficiente consiste en combinar la limpieza con el \textit{búfer de páginas}. Las páginas reemplazadas se almacenan en listas de páginas modificadas y no modificadas. Periódicamente, las páginas modificadas se escriben al almacenamiento secundario en bloques y pasan a la lista de páginas limpias, desacoplando las operaciones de limpieza y reemplazo y reduciendo el costo de las operaciones de entrada/salida.

\subsection{Control de carga}

El \textit{control de carga} determina cuántos procesos pueden permanecer simultáneamente en memoria principal, es decir, el \textit{grado de multiprogramación}.

Existe un compromiso importante:
\begin{itemize}
  \item Si hay \textit{muy pocos procesos}, el procesador puede permanecer ocioso cuando todos ellos estén bloqueados.
  \item Si hay \textit{demasiados procesos}, cada uno recibe un conjunto residente demasiado pequeño, aumentando la cantidad de fallos de página y pudiendo provocar \textit{thrashing}.
\end{itemize}

\subsubsection{Grado de multiprogramación}

El \textit{thrashing} ocurre cuando la memoria principal está sobrecargada y el sistema dedica la mayor parte del tiempo a intercambiar páginas en lugar de ejecutar procesos. Como consecuencia, la utilización del procesador disminuye drásticamente.

\begin{figure}[H]
  \centering
  % Curva clásica:
  % Eje X: grado de multiprogramación.
  % Eje Y: utilización del procesador.
  % La utilización aumenta inicialmente y luego cae debido al thrashing.
  \caption{Relación entre el grado de multiprogramación y la utilización del procesador.}
\end{figure}

Diversas estrategias permiten controlar el nivel de multiprogramación:
\begin{itemize}
  \item Los algoritmos \textit{Working Set} y \textit{PFF} ajustan dinámicamente el tamaño del conjunto residente de cada proceso, limitando indirectamente la cantidad de procesos activos.
  \item Otros criterios buscan mantener un equilibrio entre la frecuencia de fallos de página y el tiempo necesario para atenderlos, maximizando así la utilización del procesador.
  \item En algoritmos como \textit{Clock}, la velocidad con que el puntero recorre los marcos puede utilizarse como indicador de la carga del sistema: un recorrido lento sugiere memoria suficiente, mientras que uno muy rápido indica presión sobre la memoria y posible thrashing.
\end{itemize}

\subsubsection{Suspensión de procesos}

Cuando es necesario reducir el grado de multiprogramación, el sistema operativo suspende (\textit{swap out}) uno o más procesos para liberar memoria.

La selección del proceso a suspender puede basarse en distintos criterios, entre ellos:
\begin{itemize}
  \item Prioridad más baja.
  \item Proceso que produjo el fallo de página.
  \item Último proceso activado.
  \item Menor conjunto residente.
  \item Proceso de mayor tamaño.
  \item Proceso con mayor tiempo de ejecución restante.
\end{itemize}

No existe una política universalmente óptima; la elección depende de los objetivos del sistema operativo y de las características de la carga de trabajo.
